\chapter{Eigenvalues and Eigenvectors}

\section{Definitions}

Many linear operators act very simply in certain directions: a projection fixes its line and kills
the orthogonal line; a reflection fixes its line and flips the orthogonal one. Such special directions are
eigenvectors.

\begin{definition}
Let $A$ be $n\times n$. A scalar $\lambda$ is an \emph{eigenvalue} of $A$ if there is a \emph{non-zero}
vector $\vec v$ with $A\vec v=\lambda\vec v$; such a $\vec v$ is an \emph{eigenvector} of $A$ corresponding
to $\lambda$. The same terms apply to a linear operator $L$ via its standard matrix.
\end{definition}

\begin{important}[Pitfall]
$\vec 0$ is never an eigenvector (since $A\vec 0=\lambda\vec 0$ for every $\lambda$), but $0$ \emph{can} be
an eigenvalue: $0$ is an eigenvalue of $A$ iff $\Null(A)\ne\{\vec0\}$ iff $A$ is not invertible.
\end{important}

\begin{example}
For $\proj_{\vec d}$ on $\R^2$: $\vec d$ is an eigenvector with eigenvalue $1$, and any non-zero vector
orthogonal to $\vec d$ is an eigenvector with eigenvalue $0$. For $\refl_{\vec d}$ the eigenvalues are $1$
and $-1$. A rotation $R_\theta$ with $0<\theta<\pi$ has no real eigenvectors at all: no non-zero vector
keeps its direction.
\end{example}

\begin{center}
\begin{tikzpicture}[scale=0.9,>={Stealth}]
  \draw[->,gray] (-2.4,0) -- (2.8,0); \draw[->,gray] (0,-1.4) -- (0,2.8);
  \draw[thick,gray] (-0.6,-1.2) -- (1.3,2.6) node[above] {$\vec d$-line: $\lambda = 1$};
  \draw[thick,gray,dashed] (-2.2,1.1) -- (2.2,-1.1) node[below] {$\perp$ line: $\lambda = -1$};
  \draw[->,very thick,green!50!black] (0,0) -- (0.8,1.6) node[right] {$\vec v = \refl\vec v$};
  \draw[->,very thick,blue!70!black] (0,0) -- (-1.4,0.7) node[above] {$\vec w$};
  \draw[->,very thick,red!70!black] (0,0) -- (1.4,-0.7) node[above right] {$\refl\vec w = -\vec w$};
\end{tikzpicture}
\end{center}
For the reflection $\refl_{\vec d}$, vectors on the line are fixed (eigenvalue $1$) and vectors
on the perpendicular line are reversed (eigenvalue $-1$); every other vector changes direction.

\section{The characteristic polynomial}

$A\vec v=\lambda\vec v$ with $\vec v\ne\vec0$ means $(A-\lambda I)\vec v=\vec0$ has a non-trivial solution,
i.e.\ $A-\lambda I$ is not invertible, i.e.\ $\det(A-\lambda I)=0$.

\begin{definition}
The \emph{characteristic polynomial} of $A$ is $C(\lambda)=\det(A-\lambda I)$, a polynomial of degree $n$
in $\lambda$ (with leading coefficient $(-1)^n$).
\end{definition}

\begin{theorem}
$\lambda$ is an eigenvalue of $A$ iff $C(\lambda)=0$. The eigenvectors for $\lambda$ are the non-zero
vectors of $\Null(A-\lambda I)$.
\end{theorem}

\begin{definition}[Eigenspace, multiplicities]
For an eigenvalue $\lambda$, the \emph{eigenspace} is $E_\lambda=\Null(A-\lambda I)$, a subspace of
$\R^n$. The \emph{algebraic multiplicity} $a_\lambda$ is the multiplicity of $\lambda$ as a root of $C$;
the \emph{geometric multiplicity} is $g_\lambda=\dim E_\lambda$.
\end{definition}

\begin{important}[Algorithm]
(1) Compute $C(\lambda)=\det(A-\lambda I)$ and factor it. (2) For each root $\lambda$, row reduce
$A-\lambda I$ and find a basis of its null space. (Since $\lambda$ is a root, $A-\lambda I$ \emph{must}
have a non-trivial null space; if your row reduction gives only $\vec 0$, you made an arithmetic error.)
\end{important}

\begin{example}
Let $A=\begin{bmatrix}4&1\\2&3\end{bmatrix}$.
$C(\lambda)=\begin{vmatrix}4-\lambda&1\\2&3-\lambda\end{vmatrix}=(4-\lambda)(3-\lambda)-2
=\lambda^2-7\lambda+10=(\lambda-2)(\lambda-5)$.
\begin{itemize}
  \item $\lambda=2$: $A-2I=\begin{bmatrix}2&1\\2&1\end{bmatrix}\xrightarrow{R_2-R_1}\begin{bmatrix}2&1\\0&0\end{bmatrix}
        \xrightarrow{\frac12R_1}\begin{bmatrix}1&\frac12\\0&0\end{bmatrix}$, so $x_1=-\frac12t$; taking $t=-2$,
        $E_2=\Span\{(1,-2)\}$.
  \item $\lambda=5$: $A-5I=\begin{bmatrix}-1&1\\2&-2\end{bmatrix}\xrightarrow{R_2+2R_1}\begin{bmatrix}-1&1\\0&0\end{bmatrix}
        \xrightarrow{-R_1}\begin{bmatrix}1&-1\\0&0\end{bmatrix}$, so $E_5=\Span\{(1,1)\}$.
\end{itemize}
Check: $A(1,-2)=(4-2,2-6)=(2,-4)=2(1,-2)$ and $A(1,1)=(5,5)=5(1,1)$. Here $a_\lambda=g_\lambda=1$ for both.
\end{example}

Pictured: $A$ stretches the line $E_2$ by $2$ and the line $E_5$ by $5$, while a vector such
as $\vec e_1$ that lies on neither line is turned to a new direction ($A\vec e_1 = (4,2)$).
\begin{center}
\begin{tikzpicture}[scale=0.75,>={Stealth}]
  \draw[very thin,gray!30] (-1.5,-2.5) grid (4.5,3.5);
  \draw[->,gray] (-1.5,0) -- (4.8,0); \draw[->,gray] (0,-2.5) -- (0,3.8);
  \draw[thick,blue!40] (-1.2,2.4) -- (1.25,-2.5) node[below] {$E_2$};
  \draw[thick,green!40] (-1.5,-1.5) -- (3.5,3.5) node[above] {$E_5$};
  \draw[->,very thick,blue!70!black] (0,0) -- (0.5,-1) node[left] {$\vec v$};
  \draw[->,thick,blue!70!black,dashed] (0.5,-1) -- (1,-2) node[right] {$A\vec v = 2\vec v$};
  \draw[->,very thick,green!50!black] (0,0) -- (0.6,0.6) node[left] {$\vec w$};
  \draw[->,thick,green!50!black,dashed] (0.6,0.6) -- (3,3) node[above left] {$A\vec w = 5\vec w$};
  \draw[->,very thick] (0,0) -- (1,0) node[below] {$\vec e_1$};
  \draw[->,thick,red!70!black] (0,0) -- (4,2) node[right] {$A\vec e_1$};
\end{tikzpicture}
\end{center}

\begin{example}[Repeated eigenvalue]
Let $A=\begin{bmatrix}0&1&1\\1&0&1\\1&1&0\end{bmatrix}$. Expanding along the first row,
\[
C(\lambda)=\begin{vmatrix}-\lambda&1&1\\1&-\lambda&1\\1&1&-\lambda\end{vmatrix}
=-\lambda(\lambda^2-1)-1(-\lambda-1)+1(1+\lambda)
=-\lambda^3+3\lambda+2=-(\lambda+1)^2(\lambda-2).
\]
(Check: $(\lambda+1)^2(\lambda-2)=(\lambda^2+2\lambda+1)(\lambda-2)=\lambda^3-3\lambda-2$.)
\begin{itemize}
  \item $\lambda=-1$ ($a_{-1}=2$):
  $A+I=\begin{bmatrix}1&1&1\\1&1&1\\1&1&1\end{bmatrix}\xrightarrow{\substack{R_2-R_1\\R_3-R_1}}\begin{bmatrix}1&1&1\\0&0&0\\0&0&0\end{bmatrix}$.
  With $x_2=s$, $x_3=t$: $E_{-1}=\Span\{(-1,1,0),(-1,0,1)\}$, so $g_{-1}=2$.
  \item $\lambda=2$ ($a_2=1$):
  \[
  \begin{bmatrix}-2&1&1\\1&-2&1\\1&1&-2\end{bmatrix}
  \xrightarrow{R_1\leftrightarrow R_2}
  \begin{bmatrix}1&-2&1\\-2&1&1\\1&1&-2\end{bmatrix}
  \xrightarrow{\substack{R_2+2R_1\\R_3-R_1}}
  \begin{bmatrix}1&-2&1\\0&-3&3\\0&3&-3\end{bmatrix}
  \]
  \[
  \xrightarrow{\substack{R_3+R_2\\-\frac13R_2}}
  \begin{bmatrix}1&-2&1\\0&1&-1\\0&0&0\end{bmatrix}
  \xrightarrow{R_1+2R_2}
  \begin{bmatrix}1&0&-1\\0&1&-1\\0&0&0\end{bmatrix},
  \]
  so $E_2=\Span\{(1,1,1)\}$, $g_2=1$. Check: $A(1,1,1)=(2,2,2)$.
\end{itemize}
\end{example}

\begin{example}[Geometric multiplicity smaller than algebraic]
$A=\begin{bmatrix}2&1\\0&2\end{bmatrix}$ has $C(\lambda)=(2-\lambda)^2$, so $\lambda=2$ with $a_2=2$. But
$A-2I=\begin{bmatrix}0&1\\0&0\end{bmatrix}$ has rank $1$, so $E_2=\Span\{(1,0)\}$ and $g_2=1<2$.
\end{example}

\begin{example}[No real eigenvalues]
$[R_{\pi/2}]=\begin{bmatrix}0&-1\\1&0\end{bmatrix}$ has $C(\lambda)=\lambda^2+1$, with no real roots. Over
$\C$ it has eigenvalues $\pm i$. (In Math 136 we mostly work over $\R$; complex eigenvalues are
treated in Math 235.)
\end{example}

\section{Properties}

\begin{theorem}
If $A$ is upper or lower triangular, its eigenvalues are its diagonal entries (with multiplicity).
\end{theorem}
\begin{proof}
$A-\lambda I$ is triangular, so $C(\lambda)=(a_{11}-\lambda)\cdots(a_{nn}-\lambda)$.
\end{proof}

\begin{theorem}
For every eigenvalue $\lambda$, $1\le g_\lambda\le a_\lambda$.
\end{theorem}
\begin{proof}[Proof sketch]
$g_\lambda\ge1$ since an eigenvalue has an eigenvector. For the upper bound, extend a basis
$\{\vec v_1,\dots,\vec v_g\}$ of $E_\lambda$ to a basis of $\R^n$ and let $P$ have these as columns. Then
$P^{-1}AP=\begin{bmatrix}\lambda I_g & B\\ O & D\end{bmatrix}$, which has characteristic polynomial
$(\lambda-t)^g\det(D-tI)$ (in variable $t$). Similar matrices share characteristic polynomials (Chapter 10), so
$(\lambda-t)^g$ divides $C(t)$.
\end{proof}

\begin{theorem}
If $\lambda_1,\dots,\lambda_n$ are the eigenvalues of $A$ listed with algebraic multiplicity (over $\C$), then
\[
  \lambda_1+\dots+\lambda_n=\tr A \quad(\text{sum of diagonal entries}),\qquad
  \lambda_1\cdots\lambda_n=\det A .
\]
\end{theorem}
\begin{proof}[Proof sketch]
$C(\lambda)=(-1)^n(\lambda-\lambda_1)\cdots(\lambda-\lambda_n)$. Setting $\lambda=0$ gives
$\det A=C(0)=\lambda_1\cdots\lambda_n$. Comparing coefficients of $\lambda^{n-1}$ (only the product of the diagonal
entries contributes) gives the trace statement.
\end{proof}

Check on the examples: $2+5=7=4+3$ and $2\cdot5=10=12-2$; and $-1-1+2=0=\tr A$,
$(-1)(-1)(2)=2=\det A$ for the $3\times3$ example.

\begin{theorem}
If $\vec v$ is an eigenvector of $A$ for $\lambda$, then $\vec v$ is an eigenvector of $A^k$ for $\lambda^k$, and
(if $A$ is invertible) of $A^{-1}$ for $\lambda^{-1}$.
\end{theorem}
\begin{proof}
$A^2\vec v=A(\lambda\vec v)=\lambda^2\vec v$, and induct. If $A$ is invertible then $\lambda\ne0$, and applying
$A^{-1}$ to $A\vec v=\lambda\vec v$ gives $\vec v=\lambda A^{-1}\vec v$.
\end{proof}

\begin{theorem}
Eigenvectors corresponding to distinct eigenvalues are linearly independent. More generally, if
$\lambda_1,\dots,\lambda_k$ are distinct and $B_i$ is a basis of $E_{\lambda_i}$, then $B_1\cup\dots\cup B_k$ is
linearly independent.
\end{theorem}
\begin{proof}[Proof (first statement)]
Suppose not, and let $\{\vec v_1,\dots,\vec v_m\}$ be a shortest dependent list of such eigenvectors, so
$c_1\vec v_1+\dots+c_m\vec v_m=\vec0$ with all $c_i\neq0$ (by minimality). Apply $A-\lambda_mI$:
$c_1(\lambda_1-\lambda_m)\vec v_1+\dots+c_{m-1}(\lambda_{m-1}-\lambda_m)\vec v_{m-1}=\vec0$, a shorter non-trivial
relation (the coefficients are non-zero because the $\lambda_i$ are distinct). Contradiction.
\end{proof}

\section*{Summary}
\begin{itemize}
  \item $A\vec v=\lambda\vec v$, $\vec v\ne\vec 0$. Eigenvalues are roots of $C(\lambda)=\det(A-\lambda I)$;
        eigenspaces are $E_\lambda=\Null(A-\lambda I)$.
  \item $1\le g_\lambda\le a_\lambda$; the sum of eigenvalues is the trace and the product is the determinant.
  \item Triangular matrices: eigenvalues on the diagonal. $0$ is an eigenvalue iff $A$ is singular.
  \item Eigenvectors from distinct eigenvalues are linearly independent.
\end{itemize}
